\subsection{Parallel binary composition maps}

\begin{proposition}[Proved]
\label{CP2:A-four-binary-stalk-map-family}
Let $C$ have four objects $p,s,q,r$, in that order, and
hom-set cardinality matrix
\[
 \begin{pmatrix}
 1&0&0&z\\
 0&a&x&2\\
 0&0&1&2\\
 0&0&0&m
 \end{pmatrix},
 \qquad a,x,z,m\geq1,
 \qquad a=1\ \hbox{or}\ x=1.
\]
Then either every rational degree-two product on $BC$
vanishes, or there is an actual category of the same
projective-plane type, if that type exists, with
strictly fewer arrows. Consequently no such category
with at most eighteen arrows has type $\mathbb{CP}^2$.
The opposite statement also holds.
\end{proposition}
\begin{proof}
Write $N=\operatorname{End}_C(r)$,
$A=C(s,q)$, $U=C(q,r)$ and $V=C(s,r)$.
The sets $U,V$ are two-element left $N$-sets.
For each $b\in A$, composition defines the equivariant
map
\[
 \lambda_b:U\longrightarrow V,
 \qquad u\longmapsto u\circ b.
\]
The full opposite block $E=C|_{\{p,s,q\}}$ has
$H^2(BE;\mathbb Q)=0$. To see this, $p$ is isolated
within $E$. If $a=1$, its other component is a graph
with $x$ edges between two trivial corners. If $x=1$,
that component has actual terminal object $q$, regardless
of the corner at $s$. This verifies the required
vanishing even though $E$ is disconnected.

If some $\lambda_b$ is constant, its value $v\in V$
is fixed by all $n\in N$, since
$n\lambda_b(u)=\lambda_b(nu)$. The fixed-arrow
Corollary~\ref{CP2:A-fixed-arrow} applies to the cut
into $E$ and $N$ and kills all degree-two products.
We may therefore assume every $\lambda_b$ is a bijection.

If two of these bijections differ, their relative
permutation is the nonidentity permutation of $U$.
It commutes with every $N$ actor by equivariance.
The only self-maps of a two-element set commuting with
its interchange are its two permutations. Thus $N$
acts on $U$ by permutations. The action-bar projection
$U_{hN}\to BN$ induces an injection in rational
cohomology. One can verify this directly: lift each
bar simplex to the sum of its two lifts, indexed by
$U$. The permutation property makes this a chain map,
and its composite with projection is multiplication
by two. Dualizing over $\mathbb Q$ proves injectivity.

The action bar $U_{hN}$ is the bridge subspace over the
object $q$ in the two-block correspondence. Hence
pullback from $BN$ to the entire bridge is injective
as well. Together with $H^2(BE;\mathbb Q)=0$,
Mayer--Vietoris now makes every second class of $BC$
restrict to zero on both blocks. The relative
suspension in Lemma~\ref{CP2:A-two-block} again kills
all degree-two products.

It remains that all $\lambda_b$ are the same bijection.
If $x=1$, let $b$ be the unique arrow $s\to q$ and
delete $s$. In the comma category from $s$ to the
retained full subcategory, $b$ is initial: the only
other objects are the two arrows $s\to r$, and each
has exactly one morphism from $b$, by the bijection
$\lambda_b$. The corner at $q$ is trivial. For every
retained object, its identity is initial in its comma
category. Quillen's Theorem~A therefore gives an
actual homotopy equivalence after deleting $s$.
This removes its $a$ corner arrows, its arrow to $q$,
and its two arrows to $r$, hence $a+3$ arrows.

If $x\geq2$, our hypothesis gives $a=1$. Take the
support cut with source block $\{s\}$ and complementary
block $\{p,q,r\}$. Choose $b_0\in A$. On its connecting
profunctor, send every element of $A$ to $b_0$, fix
$V$ pointwise, and use the unique map on the empty
component towards $p$. This is natural. The corner
actions at $s,q$ are trivial, the $N$ action commutes
with the identity on $V$, and the only remaining
naturality equations are exactly
$\lambda_b=\lambda_{b_0}$ for all $b\in A$.
It is an idempotent proper cut retraction. By
Theorem~\ref{CP2:A-cut-retract}, if $BC\simeq\mathbb{CP}^2$
its image has that same actual homotopy type, with
$x-1\geq1$ fewer arrows.

Both reductions contradict
Theorem~\ref{CP2:A-global-lower-eighteen} when the
original category has at most eighteen arrows.
The vanishing alternatives contradict its required
nonzero rational square. Opposition gives the dual
statement.
\end{proof}

\begin{remark*}
The frozen eighteen-arrow rows $115$ and $454$ are
respectively the instances
$(a,x,z,m)=(1,2,2,7)$ and $(2,1,2,7)$.
The proof covers every choice of their corner tables,
actions and compositions; it uses no finite search
receipt for either row.
\end{remark*}
